\chapter{Estrogen Levels}

\section*{What Is This Test?}
Estrogens are a group of C18 steroid hormones that are the primary female sex hormones. The three major endogenous estrogens in order of decreasing biological potency are: estradiol (E2, the most potent and predominant estrogen in premenopausal women), estrone (E1, the predominant estrogen in postmenopausal women, produced by peripheral aromatization of androstenedione in adipose tissue), and estriol (E3, the weakest, produced in large quantities by the placenta during pregnancy). Estradiol is produced primarily by the granulosa cells of the ovarian follicle under FSH stimulation (aromatization of androstenedione from theca cells), with smaller contributions from the corpus luteum, adrenal cortex, and peripheral aromatization of testosterone in adipose tissue, liver, skin, and brain \cite{nelson2001estrogen}. In men, estradiol is produced by peripheral aromatization of testosterone (~20\%) and by the testes (Leydig cells and Sertoli cells). Approximately 60\% of circulating estradiol is bound to albumin, 38\% to sex hormone-binding globulin (SHBG), and 1--2\% is free and biologically active. Estradiol levels vary dramatically across the menstrual cycle: low in the early follicular phase (20--50 pg/mL), rising to a pre-ovulatory peak (200--400 pg/mL), falling briefly after ovulation, then rising again during the mid-luteal phase (100--200 pg/mL) from corpus luteum production. Estrogen testing may measure total estradiol, free estradiol, estrone, or estriol. The estradiol test is most commonly ordered and is used to evaluate ovarian function, menstrual disorders, infertility, menopause, pubertal development, gynecomastia in males, and estrogen-producing tumors.

\section*{Alternative Names}
Estradiol (E2); Estrone (E1); Estriol (E3); Serum Estradiol; Plasma Estradiol; Total Estrogens; Free Estradiol; Bioavailable Estradiol; Estradiol by LC-MS/MS

\section*{Principle}
Estradiol is measured using competitive immunoassays on automated analyzers. An anti-estradiol antibody competes for estradiol with labeled estradiol (chemiluminescent tracer). A displacement agent releases estradiol from SHBG and albumin. The chemiluminescent signal is inversely proportional to estradiol concentration. Common platforms: Roche Cobas ECLIA, Abbott Architect CMIA, Siemens Atellica/ADVIA Centaur CLIA, and Beckman Coulter DxI \cite{tietz2023textbook}. Estradiol immunoassays have limited sensitivity at low concentrations (<20 pg/mL), which is problematic for measuring estradiol in men, postmenopausal women, and children, where estradiol values are typically <20--40 pg/mL. Standard immunoassays also have cross-reactivity with estrone, estriol, and estrogen conjugates (2--8\%). The gold standard reference method is liquid chromatography-tandem mass spectrometry (LC-MS/MS), which offers superior sensitivity (detection limit <1--2 pg/mL) and specificity by separating estradiol from interfering estrogens and their metabolites. The 2013 Endocrine Society position statement and the 2018 CDC Hormone Standardization Program recommend LC-MS/MS for measuring low estradiol concentrations in men, postmenopausal women, and children, where immunoassays are inaccurate \cite{rosner2013challenges}. Free and bioavailable estradiol can be calculated from total estradiol, SHBG, and albumin concentrations using mass action equations.

\section*{History}
The first estrogen, estrone, was isolated from the urine of pregnant women by Adolf Butenandt in 1929 (Nobel Prize 1939). Estradiol was isolated from sow ovaries by Doisy and colleagues in 1935. The first radioimmunoassay for estradiol was developed by Guy Abraham in 1969. The complex regulation of the menstrual cycle, including the role of estradiol positive feedback in generating the LH surge, was elucidated in the 1970s by Knobil and others. The development of sensitive immunoassays in the 1980s allowed measurement of estradiol in men and postmenopausal women for the first time. The role of aromatase in peripheral estrogen production was characterized, leading to the development of aromatase inhibitors (anastrozole, letrozole) in the 1990s for breast cancer treatment. The superiority of LC-MS/MS over immunoassay for low estradiol measurement was demonstrated in multiple studies in the 2000s--2010s, leading to the 2013 Endocrine Society recommendation.

\section*{Why This Test Is Ordered}
\begin{itemize}
  \item Evaluation of ovarian function and menstrual cycle in infertility workup --- follicular phase estradiol (day 2--4) should be low (<80 pg/mL); elevated estradiol with suppressed FSH in early follicular phase suggests an estrogen-producing cyst or diminished ovarian reserve
  \item Monitoring ovarian stimulation during IVF cycles --- serial estradiol measurements track follicular development; each mature follicle (>18 mm) contributes approximately 150--250 pg/mL of estradiol
  \item Diagnosis of menopause (estradiol <20--30 pg/mL with elevated FSH >40 mIU/mL) or premature ovarian insufficiency
  \item Investigation of pubertal development --- elevated estradiol in girls <8 years suggests precocious puberty
  \item Evaluation of gynecomastia in males --- elevated estradiol or an elevated estradiol-to-testosterone ratio suggests estrogen excess contributing to breast tissue development \cite{wierman2014androgen}
  \item Diagnosis of estrogen-secreting ovarian tumors (granulosa cell tumors, theca cell tumors) or feminizing adrenal tumors --- elevated estradiol with suppressed FSH/LH
  \item Monitoring hormone replacement therapy (HRT) in postmenopausal women --- target estradiol 40--60 pg/mL for symptom relief, though levels are not routinely monitored \cite{santen2010postmenopausal}
  \item Monitoring aromatase inhibitor therapy in breast cancer --- goal is estradiol suppression to <5--10 pg/mL
  \item Investigating unexplained bone loss or fractures in men --- estradiol (not testosterone) is the primary sex steroid regulator of bone density in men
  \item Assessment of estrogen status in boys with delayed puberty or suspected hypogonadism
  \item Evaluation of postmenopausal bleeding --- elevated estradiol without adequate progesterone opposition suggests unopposed estrogen stimulation of the endometrium
\end{itemize}

\section*{Primary vs Secondary Test}
Estradiol is a primary test for evaluating ovarian function, monitoring ovarian stimulation in IVF, and evaluating pubertal development. It is a secondary (adjunctive) test in the evaluation of male hypogonadism, gynecomastia, and postmenopausal monitoring. In PCOS evaluation, estradiol is measured (with FSH) but elevated testosterone and LH are the hallmark findings.

\section*{How This Test Is Performed}
A healthcare professional draws a blood sample from a vein into a serum separator tube (gold-top) or plain red-top tube (~3--5 mL). In menstruating women, the sample MUST be drawn at a specified cycle day: day 2--4 for ovarian reserve/follicular phase evaluation; days 10--14 for peri-ovulatory estradiol peak; day 21--23 for luteal phase evaluation. Cycle day must be documented. In postmenopausal women and men, estradiol can be drawn at any time of day. Estradiol is stable at room temperature for 24 hours and refrigerated for 7 days.

\section*{What Preparation Is Needed}
No fasting is required. In menstruating women, cycle day must be documented. Oral contraceptives and estrogen therapy elevate measured estradiol or contain ethinyl estradiol which may not cross-react in estradiol assays but suppresses endogenous estradiol. High-dose biotin (>5 mg/day) must be held 48 hours before testing. Aromatase inhibitors (anastrozole, letrozole) and GnRH agonists suppress estradiol. Clomiphene elevates estradiol. SHBG levels should be measured concomitantly to calculate free and bioavailable estradiol.

\section*{Contraindications and Precautions}
\begin{itemize}
  \item Estradiol values MUST be interpreted in the context of menstrual cycle phase, patient age, sex, pregnancy status, and menopausal status. Standard immunoassays are inaccurate at low estradiol concentrations (<20--40 pg/mL) --- LC-MS/MS is the preferred method for measuring estradiol in men, postmenopausal women, and children. Cross-reactivity: estrone (E1) cross-reacts 2--5\% in estradiol immunoassays; estrogen conjugates, ethinyl estradiol (contraceptive), and phytoestrogens have variable cross-reactivity. SHBG levels profoundly affect total estradiol --- conditions that increase SHBG (estrogen therapy, pregnancy, hyperthyroidism, liver disease) increase total estradiol without increasing free estradiol; conditions that decrease SHBG (obesity, insulin resistance, androgens, hypothyroidism, nephrotic syndrome) decrease total estradiol without decreasing free estradiol. Free or bioavailable estradiol calculation is recommended for accurate assessment. Estradiol exhibits circadian variation (higher in the afternoon) and significant intra-cycle variability. Biotin supplements interfere with streptavidin-biotin assays. Hemolysis, icterus, and lipemia interfere with spectrophotometric detection.
\end{itemize}
\section*{Risks}
Minimal risk. Slight pain, bruising, or hematoma at venipuncture site. Rarely: vasovagal syncope.

\section*{What Happens After the Test}
Results available within 1--4 hours (immunoassay) or 3--7 days (LC-MS/MS). In fertility evaluation: day 3 estradiol should be <80 pg/mL; elevated levels suggest diminished ovarian reserve or a functional cyst. Serial estradiol levels during IVF guide gonadotropin dosing to achieve a gradual rise to 1000--4000 pg/mL at hCG trigger. Estradiol <20--30 pg/mL with FSH >40 mIU/mL confirms menopause. Elevated estradiol in a postmenopausal woman not on HRT requires investigation for an estrogen-secreting tumor with pelvic ultrasound and/or CT. Low estradiol in a man with osteoporosis suggests estrogen deficiency and may benefit from low-dose estrogen therapy in select cases.

\section*{Understanding Results}

\subsection*{Below Normal}
\begin{itemize}
  \item Menopause: estradiol <20--30 pg/mL with FSH >40 mIU/mL. Ovarian follicular depletion eliminates the primary source of estradiol.
  \item Premature ovarian insufficiency (POI): estradiol <20--30 pg/mL with FSH >40 mIU/mL in women <40 years. Causes: genetic (Turner syndrome, FMR1 premutation), autoimmune, iatrogenic, or idiopathic.
  \item Functional hypothalamic amenorrhea: low estradiol with low or inappropriately normal FSH/LH. Caused by excessive exercise, low body weight, eating disorders, or stress. GnRH pulsatility is suppressed.
  \item Primary or secondary hypogonadism in males: low estradiol with low testosterone. Estradiol deficiency contributes to bone loss, hot flashes, and sexual dysfunction.
  \item Aromatase deficiency: rare genetic disorder causing inability to convert androgens to estrogens. Low estradiol with elevated testosterone in both sexes. Osteoporosis, tall stature with open epiphyses in adulthood.
  \item Aromatase inhibitor therapy (anastrozole, letrozole): pharmacologic suppression of estradiol in breast cancer treatment. Target estradiol <5--10 pg/mL in postmenopausal women.
\end{itemize}

\subsection*{Above Normal}
\begin{itemize}
  \item Ovulation (pre-ovulatory peak): estradiol 200--400 pg/mL, triggering the LH surge. Sustained elevation (>200 pg/mL for >36 hours) required for positive feedback.
  \item Controlled ovarian hyperstimulation (IVF): estradiol may reach 1000--4000 pg/mL before hCG trigger. Levels >4000--5000 pg/mL are associated with increased risk of ovarian hyperstimulation syndrome (OHSS).
  \item Estrogen-secreting ovarian tumors: granulosa cell tumors (most common, produce estradiol and inhibin), theca cell tumors. Present with precocious puberty in girls, postmenopausal bleeding, or abnormal uterine bleeding in premenopausal women. Estradiol elevated, FSH/LH suppressed.
  \item Feminizing adrenal tumors: rare, produce estradiol and estrone. Associated with gynecomastia in men and precocious puberty or abnormal bleeding in females.
  \item Obesity: increased peripheral aromatization of androgens to estrogens in adipose tissue. Mildly elevated estrone and estradiol. Associated with higher breast and endometrial cancer risk in postmenopausal women.
  \item Gynecomastia in males: elevated estradiol or elevated estradiol-to-testosterone ratio. Causes: obesity, testicular tumors (Leydig cell, Sertoli cell), adrenal tumors, liver cirrhosis, hyperthyroidism, medications (spironolactone, cimetidine, marijuana), or physiologic (neonatal or pubertal).
  \item Pregnancy: estradiol rises progressively from ~200 pg/mL in first trimester to 10,000--30,000 pg/mL at term. Produced by the fetoplacental unit.
  \item Exogenous estrogen administration: hormone replacement therapy, oral contraceptives, or surreptitious use. Ethinyl estradiol may not be detected by estradiol immunoassays but suppresses FSH/LH.
  \item Liver cirrhosis: decreased hepatic metabolism and clearance of estrogens, elevated SHBG. Elevated estrone from increased peripheral aromatization. Associated with gynecomastia and spider angiomas.
\end{itemize}

\section*{Normal Results}
\begin{itemize}
  \item Adult female, follicular phase (day 2--4): 20--50 pg/mL (75--180 pmol/L); must be <80 pg/mL on day 3
  \item Adult female, pre-ovulatory peak: 150--400 pg/mL (550--1500 pmol/L)
  \item Adult female, luteal phase: 50--200 pg/mL (180--730 pmol/L)
  \item Adult female, postmenopausal: <10--30 pg/mL (<37--110 pmol/L); immunoassays overestimate at these levels --- LC-MS/MS preferred
  \item Adult male: 10--40 pg/mL (37--150 pmol/L); LC-MS/MS gives lower and more accurate values (typical range 15--35 pg/mL)
  \item Pre-pubertal children (male and female): <10 pg/mL (<37 pmol/L)
  \item Pregnancy, first trimester: 200--2,000 pg/mL (rising progressively)
  \item Pregnancy, second trimester: 2,000--10,000 pg/mL
  \item Pregnancy, term: 10,000--30,000 pg/mL
  \item Estradiol target on IVF: approximately 150--250 pg/mL per mature follicle >18 mm
  \item Note: Reference ranges vary by assay and laboratory. LC-MS/MS and immunoassay values are NOT interchangeable. Always use laboratory-specific ranges and note methodology.
\end{itemize}

\section*{Follow-up Tests}
\begin{itemize}
  \item FSH and LH: Essential co-tests for evaluating gonadal axis. The pattern of estradiol-FSH-LH classifies ovarian failure, hypogonadotropic hypogonadism, and PCOS.
  \item Progesterone (luteal phase, day 21--23): Confirms ovulation if >3--5 ng/mL. Estradiol and progesterone together confirm luteal function.
  \item SHBG: Required to calculate free and bioavailable estradiol from total estradiol. Essential when total estradiol does not match clinical picture (e.g., high total estradiol but no hyperestrogenic symptoms in a woman on oral estrogen therapy with high SHBG).
  \item Testosterone (total and free): In evaluation of PCOS, male hypogonadism, and gynecomastia \cite{bhasin2018testosterone}.
  \item Pelvic ultrasound (transvaginal): Evaluate ovaries for follicles, cysts, or tumors. Antral follicle count correlates with ovarian reserve.
  \item Pituitary MRI: If secondary hypogonadism is suspected (low FSH/LH with low estradiol).
  \item Inhibin B and AMH: Markers of ovarian granulosa cell function and ovarian reserve.
  \item Karyotype and FMR1 testing: For premature ovarian insufficiency in women <40 years.
  \item 17-hydroxyprogesterone: Exclude non-classical congenital adrenal hyperplasia as a cause of hyperandrogenism and menstrual irregularity.
  \item DHEA-S: Distinguish adrenal from ovarian source of hyperandrogenism.
\end{itemize}

\section*{References}
\bibliographystyle{unsrtnat}
\bibliography{references}

\clearpage
\section*{Regulatory Status and Authorization Date}
This chapter describes a test category, not a specific commercial product. FDA, CDSCO, EU IVDR, and MHRA approval or registration dates therefore cannot be assigned without the assay manufacturer, product name, instrument, and intended use. For laboratory-developed tests, an individual agency approval date may not exist. See the Preface for the interpretation of regulatory status.

\clearpage
